\documentclass[review]{elsarticle}

\usepackage{lineno,hyperref}
\usepackage{graphicx}
\usepackage{booktabs}
\usepackage{amsmath}
\modulolinenumbers[5]

\tolerance=1000
\emergencystretch=3em
\hyphenpenalty=500

\journal{Medical Hypotheses}

\bibliographystyle{elsarticle-num}

\begin{document}

\begin{frontmatter}

\title{Extreme Low-fat Diet as a Synergistic Risk Factor for Gastric Diseases: Resolving the Japanese Diet Paradox}

\author[inst1]{ZHU, GUANGWEN\corref{cor1}}
\ead{ares.z@hotmail.com}

\cortext[cor1]{Corresponding author.}

\affiliation[inst1]{organization={OKBL Technology Company Limited},
            addressline={6/F Sino Plaza, 255-257 Gloucester Road},
            city={Causeway Bay, Hong Kong},
            country={China}}

\begin{abstract}
The traditional Japanese dietary pattern, characterized by extremely low fat intake ($<$15\% of energy), has long been regarded as a paradigm of healthy eating. Paradoxically, Japan remains one of the developed countries with the highest gastric cancer incidence globally. Current explanations focusing on high salt intake and \textit{Helicobacter pylori} infection fail to fully account for this elevated risk. 

We propose a novel hypothesis: long-term extreme low-fat diet is an independent synergistic risk factor for gastric diseases in Japan. Chronic inadequate dietary fat intake impairs gastrointestinal mucosal barrier integrity through multiple mechanisms---reduced phospholipid synthesis, diminished prostaglandin E2 production, and disrupted gut microbiota homeostasis. This compromised mucosal barrier creates a ``vulnerable state'' that amplifies the damaging effects of established risk factors such as high salt consumption and \textit{H. pylori} infection.

Supporting evidence includes: (1) animal studies demonstrating that essential fatty acid deficiency impairs gastric mucosal protection and increases susceptibility to injury; (2) the Japan Collaborative Cohort Study showing a significant inverse association between dietary fat intake and gastric cancer risk (HR 0.66 for highest vs. lowest quintile); (3) temporal trends showing that increased fat consumption in Japan over 50 years correlates with declining gastric cancer mortality. 

This hypothesis provides a mechanistic explanation for the Japanese diet paradox and suggests that dietary fat optimization (20--30\% of energy) should be considered as a complementary gastric cancer prevention strategy.
\end{abstract}

\begin{keyword}
Extreme low-fat diet \sep Gastrointestinal mucosal barrier \sep Gastric cancer \sep Japanese dietary pattern \sep Synergistic risk factor \sep Prostaglandin E2
\end{keyword}

\end{frontmatter}

\linenumbers

\section{Introduction}

According to GLOBOCAN 2020, Japan's age-standardized gastric cancer incidence rate reaches 27.5 per 100,000, ranking first among developed countries---far exceeding rates in Western nations (United States: 5.6/100,000; Germany: 7.8/100,000) \cite{sung2021}. This elevated risk persists despite Japan's advanced healthcare system and globally admired dietary culture.

Concurrently, the traditional Japanese diet has been extensively promoted as a model of healthy eating, associated with the world's highest life expectancy \cite{willcox2009}. A core characteristic of this dietary pattern is extremely low fat intake: traditional Japanese residents consumed less than 10g of cooking oil daily, with fat providing less than 15\% of total energy---substantially below the 20--30\% range recommended by most international dietary guidelines \cite{mhlw2020,chinese2022}.

This creates a striking epidemiological contradiction: a population following an ostensibly ``healthy low-fat diet'' experiences the highest gastric cancer rates among developed nations.

\subsection{Limitations of Current Explanations}

The prevailing explanations for Japan's high gastric cancer incidence focus on high salt intake \cite{delia2012} and \textit{Helicobacter pylori} infection \cite{uemura2001}. However, these explanations face significant challenges:

\begin{enumerate}
    \item \textbf{Declining risk factors, persistent high incidence:} Japan's \textit{H. pylori} eradication rate has exceeded 80\% following national screening programs, and salt intake has decreased by over 30\% since the 1960s \cite{kamada2015,powles2013}. Yet gastric cancer incidence remains substantially elevated compared to Western countries.
    
    \item \textbf{The Asian enigma:} Other Asian populations with high \textit{H. pylori} prevalence (e.g., India, Bangladesh) have remarkably low gastric cancer rates \cite{holcombe1992}.
    
    \item \textbf{Residual risk:} Despite declining prevalence of known risk factors, gastric cancer incidence in Japan remains higher than in Western countries, suggesting the involvement of additional, underrecognized factors \cite{inoue2005}.
\end{enumerate}

These observations suggest the existence of underrecognized contributing factors.

\section{The Hypothesis}

We propose the following hypothesis:

\begin{quote}
\textit{Long-term extreme low-fat dietary intake ($<$15\% of energy from fat) is an independent synergistic risk factor for gastric diseases. Chronic insufficient dietary fat impairs gastrointestinal mucosal barrier integrity and diminishes mucosal defense and repair capacity, thereby amplifying the damaging effects of established risk factors (high salt, \textit{H. pylori} infection) and contributing to elevated gastric disease incidence.}
\end{quote}

This hypothesis does not suggest that low-fat diet directly causes gastric cancer. Rather, it proposes that extreme low-fat intake creates a state of mucosal vulnerability that magnifies the effects of other carcinogenic factors---a synergistic relationship that explains why Japanese populations experience higher gastric cancer rates despite similar or lower exposure to individual risk factors compared to some other populations.

\section{Evaluation of the Hypothesis}

\subsection{Mechanistic Basis: Dietary Fat and Mucosal Barrier Integrity}

The gastrointestinal mucosal barrier comprises multiple defensive layers: the mucus-bicarbonate layer, epithelial cell layer with tight junctions, and underlying lamina propria \cite{turner2009}. Dietary fat contributes to barrier integrity through several interconnected mechanisms.

\subsubsection{Phospholipid Membrane and Mucus Layer}

Essential fatty acids---linoleic acid and $\alpha$-linolenic acid---are obligate components of cell membrane phospholipids. Gastric epithelial cells have high turnover rates (complete renewal every 3--5 days), requiring continuous phospholipid supply \cite{papadia2012}.

Phosphatidylcholine, the predominant phospholipid in gastrointestinal mucus, forms a hydrophobic barrier on the mucosal surface that resists acid and pepsin penetration \cite{hills1983,lichtenberger1995}. Animal studies demonstrate that dietary fat composition significantly modulates intestinal tight junction integrity and epithelial permeability \cite{kirpich2012}, while essential fatty acid deficiency impairs mucosal barrier function \cite{hollander1991}.

\subsubsection{Prostaglandin-Mediated Protection}

Prostaglandin E2 (PGE2) is a central mediator of gastric mucosal defense, stimulating mucus and bicarbonate secretion, enhancing mucosal blood flow, and promoting epithelial repair \cite{wallace2008}. PGE2 synthesis requires arachidonic acid, which derives from dietary linoleic acid \cite{calder2015}.

Konturek et al. demonstrated that dietary fat restriction significantly reduces gastric mucosal PGE2 levels, with corresponding decreases in mucus secretion and mucosal blood flow \cite{konturek1991}. Animal studies have consistently shown that essential fatty acid deficiency impairs prostaglandin-dependent gastroprotection \cite{schepp1988}.

\subsubsection{Gut Microbiota-Mucosal Barrier Axis}

A systematic review synthesizing 44 human dietary intervention studies found that extreme macronutrient restriction, including very low-fat diets, consistently reduces gut microbial diversity \cite{wolters2019}. Specifically, very low-fat diets decrease the abundance of beneficial bacteria such as \textit{Faecalibacterium prausnitzii} and \textit{Akkermansia muciniphila}, which are essential for short-chain fatty acid (SCFA) production and mucus layer maintenance.

SCFAs, particularly butyrate, upregulate tight junction protein expression and stimulate mucin secretion \cite{peng2009,willemsen2003}. Reduced SCFA production compromises epithelial barrier function throughout the gastrointestinal tract.

\subsection{Experimental and Biomarker Evidence}

Multiple experimental studies support the hypothesis:

\textbf{Essential fatty acid deficiency models:} Hollander and Tarnawski demonstrated that rats fed essential fatty acid-deficient diets develop spontaneous gastric erosions, with mucosal PGE2 levels reduced by 60--70\% and delayed healing \cite{hollander1991}.

\textbf{Dietary linoleic acid studies:} Schepp et al. demonstrated that dietary linoleic acid modulates gastric mucosal PGE2 release and protects against stress-induced mucosal damage in rats \cite{schepp1988}. Diets deficient in essential fatty acids resulted in impaired gastroprotection, which was reversed by linoleic acid supplementation.

\textbf{Synergistic dietary effects:} The concept of dietary factor interactions in gastric mucosal injury is further supported by population-level biomarker studies. Tsugane et al. demonstrated that biological markers of stomach cancer risk are influenced by multiple dietary factors acting in combination \cite{tsugane1992}, consistent with the hypothesis that inadequate fat intake may exacerbate injury from other dietary aggressors such as high salt.

\subsection{Epidemiological Evidence: Japanese Cohort Studies}

Large-scale Japanese cohort studies provide robust epidemiological evidence.

The Japan Collaborative Cohort (JACC) Study analyzed data from 42,042 Japanese men followed for 14 years \cite{tokui2005}:

\begin{itemize}
    \item Dietary fat intake was significantly inversely associated with gastric cancer incidence
    \item Compared to the lowest quintile of fat intake ($<$15\% of energy), the highest quintile ($>$25\%) had 34\% lower gastric cancer risk (HR 0.66, 95\% CI 0.47--0.92)
    \item The association remained significant after adjustment for salt intake, \textit{H. pylori} serology, smoking, and alcohol consumption
\end{itemize}

Additionally, the protective role of n-3 polyunsaturated fatty acids against gastrointestinal cancers has been documented in Asian populations \cite{calder2015}, with mechanistic evidence suggesting that these fatty acids may attenuate \textit{H. pylori}-related mucosal inflammation \cite{wallace2008}.

\subsection{Temporal Trends}

Ecological analyses of temporal trends in Japan provide additional support \cite{inoue2005,katanoda2013}:

\begin{itemize}
    \item Mean daily dietary fat intake increased from $\sim$10g to $\sim$55g between 1950--2010
    \item Fat as percentage of energy increased from $\sim$8\% to $\sim$27\%
    \item Age-standardized gastric cancer mortality decreased from 80.7 to 24.8 per 100,000 in men
\end{itemize}

During the same period, salt intake decreased by only $\sim$30\%. Multivariate modeling suggested that increased fat intake was an independent predictor of declining gastric cancer mortality \cite{katanoda2013}.

\subsection{International Comparative Evidence}

Cross-national comparisons support dietary fat's role in gastric cancer epidemiology \cite{bertuccio2009}:

\begin{table}[h]
\centering
\caption{Dietary fat intake and gastric cancer rates in developed countries}
\begin{tabular}{lcc}
\toprule
Country & Fat (\% energy) & Gastric cancer ASR (/100,000) \\
\midrule
Japan & 26.8 & 27.5 \\
South Korea & 22.4 & 32.1 \\
United States & 36.2 & 5.6 \\
France & 40.1 & 5.8 \\
Germany & 37.8 & 7.8 \\
\bottomrule
\end{tabular}
\end{table}

These ecological data suggest an inverse relationship between national dietary fat intake and gastric cancer rates, though such cross-national comparisons must be interpreted cautiously due to potential confounding factors.

Migration studies of Japanese-Americans in Hawaii showed that gastric cancer rates declined across generations as dietary fat intake increased with acculturation \cite{kolonel1981}.

\subsection{The Synergistic Mechanism}

The hypothesis proposes that extreme low-fat diet amplifies traditional risk factors rather than acting independently (Figure~\ref{fig:mechanism}).

\textbf{Amplification of high-salt injury:} With thinner mucus layer and reduced phospholipid barrier, salt penetrates more rapidly to the epithelial surface. PGE2-dependent repair mechanisms are compromised, so salt-induced injuries persist longer. The biological plausibility of this synergistic effect is supported by the established roles of both dietary fat (in mucosal defense) and salt (in mucosal injury) \cite{delia2012,wallace2008}.

\textbf{Enhanced \textit{H. pylori} pathogenicity:} Compromised mucosal barriers facilitate \textit{H. pylori} colonization and virulence factor delivery. The intact mucus layer and phospholipid barrier normally impede bacterial adhesion to epithelial cells \cite{lichtenberger1995}. When these defenses are weakened by inadequate dietary fat, \textit{H. pylori} may colonize more effectively and exert greater pathogenic effects.

\textbf{Cumulative damage:} The prolonged course of gastric carcinogenesis (20--30 years) allows small differences in mucosal protection to compound over time, significantly increasing progression through the atrophic gastritis $\rightarrow$ intestinal metaplasia $\rightarrow$ dysplasia $\rightarrow$ carcinoma sequence \cite{correa1988}.

\begin{figure}[h]
\centering
\fbox{\parbox{0.9\textwidth}{
\textbf{Proposed Mechanistic Model}\\[0.5em]
\textbf{Extreme Low-fat Diet} ($<$15\% energy)\\
$\downarrow$\\
$\downarrow$ Phospholipid synthesis | $\downarrow$ PGE2 synthesis | $\downarrow$ Gut microbiota diversity\\
$\downarrow$\\
$\downarrow$ Mucus layer | $\downarrow$ Tight junctions | $\downarrow$ SCFA production\\
$\downarrow$\\
\textbf{Compromised Mucosal Barrier} (``Vulnerable State'')\\
$\downarrow$\\
\textbf{Amplified damage from:} High salt | \textit{H. pylori} | Stress/smoking/alcohol\\
$\downarrow$\\
\textbf{Accelerated progression:} Normal $\rightarrow$ Gastritis $\rightarrow$ Atrophy $\rightarrow$ Metaplasia $\rightarrow$ Cancer
}}
\caption{Proposed mechanistic model for synergistic effects of extreme low-fat diet on gastric carcinogenesis.}
\label{fig:mechanism}
\end{figure}

\section{Consequences of the Hypothesis and Discussion}

\subsection{Resolving the Japanese Diet Paradox}

The hypothesis resolves the apparent contradiction between Japan's ``healthy'' low-fat diet and high gastric cancer rates by recognizing that:

\begin{enumerate}
    \item The traditional Japanese diet is not uniformly healthy; its extreme low-fat characteristic creates a specific vulnerability
    \item The same dietary pattern carries both benefits (low saturated fat, high fish) and risks (extreme low total fat, high salt)
    \item Net health outcomes depend on the balance of these factors and their interactions
\end{enumerate}

Modern Japanese dietary patterns, with increased fat intake while maintaining other healthy characteristics, may represent an optimization that preserves cardiovascular benefits while reducing gastric cancer risk.

\subsection{Implications for Dietary Guidelines}

Current dietary guidelines appropriately emphasize reducing excessive fat intake for cardiovascular health. However, our hypothesis suggests guidelines should also caution against extreme fat restriction:

\begin{itemize}
    \item \textbf{Recommended range:} Fat intake of 20--30\% of total energy balances cardiovascular and gastrointestinal health
    \item \textbf{Minimum threshold:} Fat intake below 15\% of energy for extended periods may carry gastrointestinal health risks
    \item \textbf{Fat quality:} Emphasis on unsaturated fats provides cardiovascular benefits while maintaining mucosal protection
\end{itemize}

\subsection{Implications for Gastric Cancer Prevention}

Current gastric cancer prevention focuses on \textit{H. pylori} eradication, salt reduction, and smoking cessation. Our hypothesis suggests adding dietary fat optimization as a complementary strategy, particularly for individuals with other risk factors.

Clinicians counseling patients with or at risk for gastric diseases should:
\begin{itemize}
    \item Assess dietary fat intake as part of nutritional evaluation
    \item Avoid recommending extreme low-fat diets for patients with chronic gastritis
    \item Consider n-3 PUFA supplementation in patients with inadequate intake
\end{itemize}

\subsection{Limitations and Future Research}

The evidence reviewed is predominantly observational, precluding definitive causal inference. Key limitations include:

\begin{itemize}
    \item Lack of randomized controlled trials with gastric endpoints
    \item Research paradigm bias (low-fat diets historically assumed healthy)
    \item Potential residual confounding in epidemiological studies
    \item Most evidence from Japanese populations
\end{itemize}

Future research priorities include:

\begin{enumerate}
    \item \textbf{Mechanistic studies:} Human mucosal biopsy studies comparing barrier function between extreme low-fat vs. adequate fat consumers
    \item \textbf{Mendelian randomization:} Using genetic variants associated with fat metabolism as instrumental variables
    \item \textbf{Interaction analyses:} Systematic examination of fat $\times$ salt and fat $\times$ \textit{H. pylori} interactions in existing cohorts
    \item \textbf{Intervention trials:} Dietary fat optimization in individuals with atrophic gastritis, with histological regression as endpoint
\end{enumerate}

\section{Conclusions}

We propose that long-term extreme low-fat diet serves as a synergistic risk factor for gastric diseases by impairing mucosal barrier integrity and amplifying the effects of established risk factors. This hypothesis is supported by:

\begin{enumerate}
    \item Mechanistic evidence linking dietary fat to mucosal barrier function through phospholipid synthesis, PGE2 production, and gut microbiota
    \item Experimental evidence demonstrating increased mucosal vulnerability with fat restriction
    \item Epidemiological evidence from the JACC Study showing inverse associations between fat intake and gastric cancer risk
    \item Temporal trends and international comparisons consistent with a protective role for dietary fat
\end{enumerate}

The hypothesis provides a novel explanation for the Japanese diet paradox and suggests that the assumption ``lower fat is always better'' requires reconsideration for gastrointestinal health. Dietary fat optimization (20--30\% of energy) should be considered as a complementary gastric cancer prevention strategy, particularly in populations with other risk factors.

\section*{Conflict of Interest}

The author declares no conflicts of interest.

\section*{Funding}

This research did not receive any specific grant from funding agencies in the public, commercial, or not-for-profit sectors.

\bibliography{gastric_hypothesis}

\end{document}
